\chapter{Further constructions in the source library}
\label{app:construction-details}

The preceding chapters follow the manuscript arguments. This appendix
records related auxiliary constructions in the source library. Their
hypotheses concern the same characters, groups, coefficient fields and
actions as their conclusions. They do not provide instances of the full
application hypotheses.

\section{Central kernels and local ordinary characters}

Let $X$ be a finite perfect group, let $\ell$ be prime, and suppose that
$\ell\nmid |Z(X)|$. Fix an algebraically closed modular field $k$ of characteristic $\ell$,
an ordinary field $K$ of characteristic zero and a root correspondence
$\iota$. For a representation $\rho_\varphi$ affording an irreducible
Brauer character $\varphi$, put
\[
 Z_\varphi=Z(X)\cap\ker\rho_\varphi,
 \qquad \overline X_\varphi=X/Z_\varphi.
\]
Factoring $\rho_\varphi$ through this quotient preserves irreducibility.
With the induced root convention, inflation recovers $\varphi$.
Perfectness shows that the centre of the quotient is the image of
$Z(X)$: commutators with an element central modulo $Z_\varphi$ define
a homomorphism from $X$ to the central subgroup $Z_\varphi$, and that
homomorphism is trivial. The quotient character is therefore faithful
on the quotient centre.

For a weight $(Q,\theta)$ in the block of $\varphi$, the local reduction
and block compatibility hypotheses identify the scalar character of
its local Brauer character with that of $\varphi$ on $Z(X)$.
Consequently $Z_\varphi$ is also in the local modular kernel. Root
agreement on these central elements transfers this equality to the
ordinary character $\theta$. Under the specified normaliser quotient
and radicality hypotheses, ordinary descent constructs
$\overline\theta$ with
\[
 \theta=\overline\theta\circ q_Q,
 \qquad
 q_Q:N_X(Q)/Q\longrightarrow
 N_{\overline X_\varphi}(\overline Q)/\overline Q.
\]
The map $q_Q$ is the map induced by this central quotient.
Its surjectivity gives uniqueness of $\overline\theta$.
Block descent additionally uses the complete quotient block
decompositions and preservation of the relevant primitive idempotents;
the character factorisation alone does not supply those facts.

\begin{samepage}
\begin{implementation}
The modules
\module{ModularRep.PaperProofs.SporadicFi24P3Definition44NamedCarrierFullCenterScalar}
and
\module{ModularRep.PaperProofs.SporadicFi24P3Definition44NamedCarrierOwnCentralQuotient}
contain the scalar and quotient constructions. The former proves
\lean{full_center_scalar_eq}; the latter constructs
\lean{ownCentralQuotientBrauerSource}. The hypotheses on local block
operations, quotient normalisers and coefficient conventions are retained.
\end{implementation}
\end{samepage}

An equivariant block-preserving bijection
$\Omega:\IBr_\iota(X)\simeq\mathcal W(X)$ also determines ordinary
local maps. For each radical subgroup $Q$, take the Brauer fibre
\[
 \{\varphi:\rad(\Omega\varphi)=[Q]\}.
\]
The map $\theta\mapsto[(Q,\theta)]$ identifies the defect zero
characters of $N_X(Q)/Q$ with weight classes whose radical subgroup
has conjugacy class $[Q]$. Its inverse defines $\Omega_Q$ by
\[
 [(Q,\Omega_Q(\varphi))]=\Omega(\varphi).
\]
Injectivity gives the covariance of these local maps under
automorphisms stabilising $Q$. The partition retains empty fibres.
Its nonempty parts are distinguished by their radical classes;
equality of two empty parts does not identify their labels.

\section{Extensions and compatible roots}

The local extension construction for $\varphi$ uses a finite ambient
group $A$ with a normal subgroup $B\cong\overline X_\varphi$ and a map
$f:X\to A$ with kernel $Z_\varphi$. For $D=N_A(f(Q))$, the specified
local map $f_N:N_X(Q)\to D$ satisfies
\[
 \iota_D f_N=f\iota_{N_X(Q)},
 \qquad \im(f_N)=B\cap D.
\]
The ambient hypotheses include $C_A(B)=Z(A)$ and
$\ell\nmid |Z(A)|$. The quotient $A/Z(A)$ identifies with the stabiliser of the original
character, using the opposite automorphism group to implement the
manuscript's right action. The witnesses include global and local
Brauer extensions, their irreducible restrictions for every
$B\leq J\leq A$, and the required block induction relations.

The raw witness permits separate root correspondences. The module
\module{ModularRep.PaperProofs.SpathRootCoherence} adds the agreement
conditions needed by the full inductive certificate.
For \lean{OriginalPacketRoots}, the original and quotient conventions
agree on the quotient roots; the global extension convention agrees
with the quotient and local extension conventions, and with the
conventions on $J$ and $D\cap J$ for every intermediate $J$.
When the original group embeds in the extension group,
\lean{ActualPacketRoots} additionally requires agreement with its
original roots. Factoring through the central kernel is what permits
the first construction to use the weaker condition.

Let $H$ be a finite group and $N\lhd H$ with $[H:N]\leq 2$.
A finite dimensional irreducible representation $\rho$ of $N$ over an
algebraically closed field, invariant under conjugation by $H$, extends
to $H$ on the same vector space. In the index two case, choose
$t\in H\setminus N$ and an invertible
intertwiner $T$. Schur's lemma makes
$\rho(t^2)^{-1}T^2$ a nonzero scalar $c$. Rescaling $T$ by a square
root of $c^{-1}$ gives an operator $U$ satisfying
$U^2=\rho(t^2)$; together with $\rho$, this defines the representation
on both cosets. The construction also applies to a nonsplit extension.
It is implemented in
\module{ModularRep.PaperProofs.SporadicFi24P3Definition44NamedCarrierIndexTwoExtension}.

These extension constructions do not turn a numerical bijection into
the full inductive condition without the remaining block and root
hypotheses. In the triple-cover applications, compatible replacement
witnesses are supplied for the same numerical correspondence, as
described in the sporadic chapter. Their existence remains an assumption.

\section{Selected Type B constructions}

The rank-three Morita construction keeps one Levi character, its
representation, its actual field stabiliser and its chosen Jordan
image throughout. The base functor is the balanced tensor product
with the specified block bimodule. The upper functor acts on the
extensions by the same field stabiliser. Their restriction square
gives
\[
 \operatorname{Res}_G\widehat{\mathcal F}
 \cong\mathcal F\operatorname{Res}_L.
\]
Applying the upper functor transfers a Levi extension to the ambient
group. Conversely, essential surjectivity of the upper equivalence
and full faithfulness of the base equivalence recover an extension
of the original Levi representation. The stated equality of tensor
characters with the Jordan map identifies the resulting character.
These are hypotheses and deductions for the selected bimodule and
actions, rather than an interpretation of an arbitrary abstract
Morita equivalence.

\begin{implementation}
The selected result is
\lean{ModularRep.PaperProofs.TypeBRankThreeMoritaSelectedApplication.actual_selected_morita_extension}.
The separate module
\module{ModularRep.PaperProofs.TypeBRankThreeSameFieldAutomorphism}
retains the identification between the relevant outer action and the
normalised field action on this same union of blocks. The field
stabiliser is not identified with a full automorphism stabiliser by
the Morita equivalence alone.
\end{implementation}

For a finite direct product $G=\prod_iG_i$ of finite groups, a radical $2$-subgroup $R$ is
the product of its coordinate images $R_i$. The argument uses
radicality, the normaliser condition for finite $2$-groups and the
product formula for the $2$-core. The actual coordinate maps then give
\[
 N_G(R)/R\cong\prod_iN_{G_i}(R_i)/R_i.
\]
The supplied ordinary tensor realisations, their unique factorisation
and the defect zero criterion identify the defect zero characters on
these quotients. Passing to conjugacy
classes gives the product correspondence on weights. To retain blocks,
the proof separately uses the ordinary block product identities and
checks the central character equation along the product normaliser
inclusion. The resulting
\lean{ModularRep.PaperProofs.TypeBRankThreeProductRawWeightBlock.rawWeightBlock_primitive}
therefore identifies the ambient primitive block induced from the
actual ordinary local blocks.
\begin{implementation}
The three Type B modules named in this section are not imported by any
other module in the current formalisation. They record auxiliary
constructions; this section does not assert their integration into the
manuscript deductions.
\end{implementation}
